% Options for packages loaded elsewhere
\PassOptionsToPackage{unicode}{hyperref}
\PassOptionsToPackage{hyphens}{url}
\documentclass[
]{article}
\usepackage{xcolor}
\usepackage{amsmath,amssymb}
\setcounter{secnumdepth}{-\maxdimen} % remove section numbering
\usepackage{iftex}
\ifPDFTeX
  \usepackage[T1]{fontenc}
  \usepackage[utf8]{inputenc}
  \usepackage{textcomp} % provide euro and other symbols
\else % if luatex or xetex
  \usepackage{unicode-math} % this also loads fontspec
  \defaultfontfeatures{Scale=MatchLowercase}
  \defaultfontfeatures[\rmfamily]{Ligatures=TeX,Scale=1}
\fi
\usepackage{lmodern}
\ifPDFTeX\else
  % xetex/luatex font selection
\fi
% Use upquote if available, for straight quotes in verbatim environments
\IfFileExists{upquote.sty}{\usepackage{upquote}}{}
\IfFileExists{microtype.sty}{% use microtype if available
  \usepackage[]{microtype}
  \UseMicrotypeSet[protrusion]{basicmath} % disable protrusion for tt fonts
}{}
\makeatletter
\@ifundefined{KOMAClassName}{% if non-KOMA class
  \IfFileExists{parskip.sty}{%
    \usepackage{parskip}
  }{% else
    \setlength{\parindent}{0pt}
    \setlength{\parskip}{6pt plus 2pt minus 1pt}}
}{% if KOMA class
  \KOMAoptions{parskip=half}}
\makeatother
\setlength{\emergencystretch}{3em} % prevent overfull lines
\providecommand{\tightlist}{%
  \setlength{\itemsep}{0pt}\setlength{\parskip}{0pt}}
\usepackage{bookmark}
\IfFileExists{xurl.sty}{\usepackage{xurl}}{} % add URL line breaks if available
\urlstyle{same}
\hypersetup{
  pdftitle={Homework 11 Written},
  pdfauthor={Ewan Pedersen},
  hidelinks,
  pdfcreator={LaTeX via pandoc}}

\title{Homework 11 Written}
\author{Ewan Pedersen}
\date{2025-11-21}

\begin{document}
\maketitle
\begin{abstract}
Show the steps of deriving your answers. Points will be deducted for
answers without adequate steps discussed. Submit your homework via
Brightspace as one PDF or Word document.
\end{abstract}

\begin{center}\rule{0.5\linewidth}{0.5pt}\end{center}

\subsection{Maximum Bipartite Matching: Maxflow
Solution}\label{maximum-bipartite-matching-maxflow-solution}

Given the bipartite graph(\emph{Figure 1}) shown below that we used in
class, solve the problem of finding a maximum matching by reducing the
problem to a maximum flow finding problem. Specifically, show the
resulting flow network and residual graph after finding the maximum flow
(using the basic Ford-Fulkerson algorithm). The flow network should show
on every edge the flow value assigned and the capacity (e.g., f/c) and
the residual graph should show on every edge the residual capacity.
Additionally, write the maximum flow value and the minimum cut (e.g.,
(\(\{ s,a,b,c \}, \{ d,e,t \}\)) resulting from the algorithm. Note the
flow values on the individual edges of the flow network may vary a bit
depending on the augmenting paths found randomly during the execution of
the algorithm, but the maximum flow value is certainly the same and the
minimum cut found is observably the same, too.

!{[}{[}helpme.png{]}{]}

\subsubsection{Initial Bipartite Graph
Edges}\label{initial-bipartite-graph-edges}

\[
\begin{align}
E = \{(1, 1'), (1, 2'), (2, 2'), (3, 1'), (3, 2'), \\
(3, 3'), (4, 3'), (4, 4'), (5, 4'), (5, 5') \} 
\end{align}
\]

\subsubsection{Construction of Flow
Network}\label{construction-of-flow-network}

We add the source \(s\) connected to all the leftside vertices
\(\{ 1,2,3,4,5 \}\) with a capacity of 1. Then add a sink \(t\)
connected from all the rightside vertices \(\{ 1',2',3',4',5' \}\) with
a capacity of 1. Thus all bipartite graph edges have a capacity of 1.

\subsubsection{Execution of Ford-Fulkerson for Augmenting
Paths}\label{execution-of-ford-fulkerson-for-augmenting-paths}

\[
\begin{align}
 \text{Path 1:} \quad & s \to 1 \to 1' \to t \quad \text{(flow = 1)} \\
 \text{Path 2:} \quad & s \to 2 \to 2' \to t \quad \text{(flow = 1)} \\
 \text{Path 3:} \quad & s \to 3 \to 3' \to t \quad \text{(flow = 1)} \\
 \text{Path 4:} \quad & s \to 4 \to 4' \to t \quad \text{(flow = 1)} \\
 \text{Path 5:} \quad & s \to 5 \to 5' \to t \quad \text{(flow = 1)} 
\end{align}
\]

Thus there exists no other augmenting pairs.

\subsubsection{Final Resulting Network
Flow}\label{final-resulting-network-flow}

Source edges:

\[
\begin{align} s \to 1 &: 1/1 \\ s \to 2 &: 1/1 \\ s \to 3 &: 1/1 \\ s \to 4 &: 1/1 \\ s \to 5 &: 1/1 \end{align}
\]

Bipartite Edges:

\[
\begin{align} 1 \to 1' &: 1/1 \quad \text{ (in matching)} \\
1 \to 2' &: 0/1  \\
2 \to 2' &: 1/1 \quad \text{ (in matching)} \\
3 \to 1' &: 0/1 \\
 3 \to 2' &: 0/1 \\
3 \to 3' &: 1/1 \quad \text{ (in matching)} \\
 4 \to 3' &: 0/1 \\
4 \to 4' &: 1/1 \quad \text{ (in matching)} \\
 5 \to 4' &: 0/1 \\
5 \to 5' &: 1/1 \quad \text{ (in matching)} \end{align}
\]

Sink edges:

\[
\begin{align} 1' \to t &: 1/1 \\
 2' \to t &: 1/1 \\
3' \to t &: 1/1 \\
 4' \to t &: 1/1 \\
5' \to t &: 1/1
\end{align}
\]

\subsubsection{Residual Graph with only Residual
Capacities}\label{residual-graph-with-only-residual-capacities}

Forward Residual Edges:

\[
\begin{align} s \to \{1, 2, 3, 4, 5\} &: 0 \text{ (all saturated)}  \\
1 \to 2' &: 1 \\
 3 \to 1' &: 1 \\
 3 \to 2' &: 1 \\
 4 \to 3' &: 1 \\
 5 \to 4' &: 1 \\
 \{1', 2', 3', 4', 5'\} \to t &: 0 \text{ (all saturated)} \end{align}
\]

Backward Residual Edges:

\[
\begin{align} \{1, 2, 3, 4, 5\} \to s &: 1 \text{ each} \\
 1' \to 1 &: 1 \\
 2' \to 2 &: 1 \\
 3' \to 3 &: 1 \\
 4' \to 4 &: 1 \\
 5' \to 5 &: 1 \\
 t \to \{1', 2', 3', 4', 5'\} &: 1 \text{ each} \end{align}
\]

\textbf{Results:}

Max flow: 5

Max matching for \(M\): \(\{ (1,1'), (2,2'), (3,3'), (4,4'), (5,5') \}\)
(with a perfect match)

Min Cut: \((\{ s \}, \{ 1,2,3,4,5,1',2',3',4',5',t \})\)

Thus, the cut with max capacity of 5 confirms the max flow flow/min cut
theorem, where \(Max\ Flow = Min\ Cut\ Capacity =5\)

\begin{center}\rule{0.5\linewidth}{0.5pt}\end{center}

\subsection{Software Porting}\label{software-porting}

\begin{quote}
Do the textbook exercise 29 in Chapter 7. Show how to solve this problem
by reformulating it to the problem of finding a minimum cut from a flow
network built from the set \(V\) of \(n\) applications
\(v_{1},v_{2},\dots,\) and \(v_{n}\) in a manner similar to the approach
used in the image segmentation problem (see Section 7.10). State clearly
how a flow network is constructed for a given instance of the software
porting problem.

Hint. One plausible way to construct a flow network from a software
porting problem can be to let the source node s represent the new system
and let the sink node \(t\) represent the application 1 which cannot
move to the new system. Then, there's a parallel between the software
porting problem and the image segmentation problem:

\begin{itemize}
\tightlist
\item
  The porting benefit \(b_{i}\) of an application \(i\) in the software
  porting problem parallels the foreground likelihood \(a_{i}\) of a
  pixel \(i\) in the image segmentation problem.
\item
  The separation expense \(x_{ij}\) between two applications \(i\) and
  \(j\) in the software porting problem (when only one of them is ported
  to the new system) parallels the separation penalty pij between two
  neighboring pixels \(i\) and \(j\) in the image segmentation problem
  (when one is in the foreground and the other is in the background).
\end{itemize}

\textbf{Chapter 7 Exercise 29}:

Some of your friends have recently graduated and started a small com-
pany, which they are currently running out of their parents' garages in
Santa Clara. They're in the process of porting all their software from
an old system to a new, revved-up system; and they're facing the
following problem.

They have a collection of \(n\) soft, rare applications,
\(\{ 1,2,\dots,n \}\), running on their old system; and they'd like to
port some of these to the new system. If they move application \(i\) to
the new system, they expect a net (monetary) benefit of \(b_{i}\geq 0\).
The different software applications interact with one another; if
applications \(i\) and \(j\) have extensive interaction, then the
company, will incur an expense if they move one of \(i\) or \(j\) to the
new system but not both; let's denote this expense by \(x_{ij}\geq 0\).

So, if the situation were really this simple, your friends would just
port all \(n\) applications, achieving a total benefit of
\(\sum_{i}b_{i}\). Unfortunately, there's a problem \ldots{}

Due to small but fundamental incompatibilities between the two systems,
there's no way to port application 1 to the new system; it will have to
remain on the old system. Nevertheless, it might still pay off to port
some of the other applications, accruing the associated benefit and
Incurring the expense of the interaction between applications on
different systems.

So this is the question they pose to you: Which of the remaining
applications, ff any, should be moved? Give a polynomial-time algorithm
to find a set \(S \subseteq \{ 2,3,\dots,n \}\)for which the sum of the
benefits minus the expenses of moving the applications in \(S\) to the
new system is maximized.
\end{quote}

\subsubsection{Min Cut Solution}\label{min-cut-solution}

We know that \(n\) applications \(\{ 1,2,\dots ,n \}\), that application
1 cannot be ported as must stay on the old system, the benefits
\(b_{i} \geq 0\) for any porting application \(i\), and seperation
expenses \(x_{ij}\geq 0\) if \(i\) and \(j\) are on different systems.

We first seek to find \(S \subseteq \{ 2,3,\dots, n \}\) to maximize:

\[
\begin{align} \text{Profit}(S) = \sum_{i \in S} b_i - \sum_{\substack{i \in S, j \notin S \\ \text{or } j \in S, i \notin S}} x_{ij} \end{align}
\]

We then define a new flow network construction where:

\begin{itemize}
\tightlist
\item
  \(s\): source representing new system
\item
  \(t\) sink representing app 1 and old system
\item
  \(i \in \{ 2,3,\dots,n \}\): note for each application
\end{itemize}

The edges and capacities can then be defined as:

\[
\begin{align} \text{Type 1: } & s \to i \text{ for each } i \in \{2, 3, \dots, n\}, \quad \text{capacity } = b_i \\
 \text{Type 2: } & i \to t \text{ for each } i \in \{2, 3, \dots, n\}, \quad \text{capacity } = x_{i1} \\
 \text{Type 3: } & i \leftrightarrow j \text{ for each pair } i, j \in \{2, 3, \dots, n\}, \quad \text{capacity } = x_{ij} \end{align}
\]

For the type 3 edges, we have to use undirected edges with a capacity of
\(x_{ij}\), which can be implemented as two directed edges: \(i\to j\)
witha capacity of \(x_{ij}\) and \(j\to i\) with capacity \(x_{ij}\).

We can then define a cut \((A,B)\) as a separation of \(s\) from \(t\),
where:

\begin{itemize}
\tightlist
\item
  \(A\): \(\{ s \} \cup \{ \text{applications ported to new system} \}\)
\item
  \(B\):
  \(\{ t \} \cup \{ \text{applications staying on old system with an app 1} \}\)
\end{itemize}

We then let \(S=A \cap \{ 2,3,\dots,n \}\) be the product set of ported
apps.

Calculating cut capacity with edges crossing \(A\) to \(B\):

\[
\begin{align} \text{Cut capacity} &= \sum_{j \in B \setminus \{t\}} c(s, j) + \sum_{i \in S} c(i, t) + \sum_{\substack{i \in S, j \in B \setminus \{t\}}} c(i, j) \\
 &= \sum_{j \notin S, j \geq 2} b_j + \sum_{i \in S} x_{i1} + \sum_{\substack{i \in S, j \notin S \\
 i,j \geq 2}} x_{ij} \end{align}
\]

We then connect the cut to objective by letting \(B_{\text{total}}\) be
the total benefit for all movable applications \(\sum_{i=2}^n\), writing
this as:

\[
\begin{align} \text{Profit}(S) &= \sum_{i \in S} b_i - \left[\sum_{i \in S} x_{i1} + \sum_{\substack{i \in S, j \notin S \\ i,j \geq 2}} x_{ij}\right] \\
 &= \left(B_{\text{total}} - \sum_{j \notin S, j \geq 2} b_j\right) - \sum_{i \in S} x_{i1} - \sum_{\substack{i \in S, j \notin S \\ i,j \geq 2}} x_{ij} \\
 &= B_{\text{total}} - \left[\sum_{j \notin S, j \geq 2} b_j + \sum_{i \in S} x_{i1} + \sum_{\substack{i \in S, j \notin S \\ i,j \geq 2}} x_{ij}\right] \\ &= B_{\text{total}} - \text{Cut capacity}(A, B) \end{align}
\]

Yielding a final alogorithm of:

\[
\begin{align} &\textbf{Algorithm: Software Porting via Min-Cut} \\ &\text{1. Construct flow network as described above} \\ &\text{2. Find minimum s-t cut using max-flow algorithm (e.g., Ford-Fulkerson)} \\ &\text{3. Let } S = \{i \in \{2, \dots, n\} : i \text{ is on s-side of min-cut}\} \\ &\text{4. Return } S \text{ as the set of applications to port} \end{align}
\]

With a time complexity of \(O(mn^2)\) ( or \(O(n^3)\) depending on
maxflow algorithm), where m is num edges in the network.

\end{document}
